% !TeX root = ../main.tex

\chapter{結論與未來展望}

\section{研究總結}

本研究成功開發並驗證了一套結合樂高（LEGO Technic）精密結構件與運算成像（Computational Imaging）技術的「多功能螢光倒立式顯微鏡」。相較於傳統昂貴且封閉的商用成像系統，以及在剛性與行程上存有缺陷的純 3D 列印開源方案，本研究所提出之設計典範在以下三個層次均展現出顯著的科學與工程貢獻。

在硬體結構方面，透過利用樂高積木優異的射出成型精度（公差 $<10\,\mu$m）與標準化 8.0 mm 網格幾何特性，本研究成功建構出一套具備亞微米級定位精度且穩定性極高的龍門式掃描載台（Gantry-style Stage）。這不僅突破了傳統開源顯微鏡行程受限的瓶頸，更提供了一套高度可重構（Reconfigurable）的底層架構，讓後續研究者能夠在數分鐘內完成光學架構的改裝，徹底解放了實驗設計的自由度。

在軟體與演算法層次，本研究針對低成本 CMOS 感測器的物理缺陷，實作了基於愛因斯坦求和約定的空間變異張量校正（Spatially Varying Tensor Unmixing），大幅提升了顯微影像的色彩保真度與空間均勻性。此外，自動對焦狀態機（Autofocus State Machine）的三階段搜尋邏輯，使系統能在不到 14 秒內完成長達近一毫米行程的精確鎖焦，充分驗證了演算法在實際應用中的高度魯棒性（Robustness）。

在超解析度成像層次，圖案化照明傅立葉疊層成像（piFP）技術的導入初步證實了在不更換物鏡的前提下，透過演算法與隨機散斑照明的結合，能夠突破物鏡的物理繞射極限，達成超解析度成像的目標。

\section{系統限制與討論}

儘管本系統在初步實驗中展現出令人滿意的效能，但在追求極致成像品質的過程中，仍觀察到以下待優化的物理與計算限制，有待後續研究加以改進。

在光學遮蔽方面，樂高積木的接合處並非完全氣密。在進行長時間、極弱螢光訊號的曝光過程中，環境光線仍可能透過細微間隙滲入光路，增加影像的背景底噪。此外，ABS 材質在環境溫度波動較大時仍存在微小的熱膨脹，在追求奈米等級定位精度的應用場景下需加以考量。

在計算效能方面，piFP 疊代重建演算法目前在嵌入式平台（Raspberry Pi 4）上的運算時間仍偏長，尚無法達成影像擷取後的即時（Real-time）顯示，這對於需要快速搜尋目標樣本區域的實驗場景而言仍有改進空間。

\section{未來改進與擴充方向}

針對本研究所發現之挑戰，後續的研究工作可朝以下方向深化，以進一步提升系統的實用性與科學價值。

\subsection{硬體結構之遮光強化}

未來計畫設計一套專屬的黑色遮光罩系統，在關鍵光路接合處採用不透光的橡膠墊片，以徹底排除雜散光對螢光訊號的干擾。同時，可嘗試在樂高結構的關鍵承重節點內崁入金屬支撐桿，進一步提升龍門架構在搭載重型高倍率物鏡時的動態剛性。

\subsection{基於邊緣運算的即時重建加速}

為了提升 piFP 的易用性，未來可將重建演算法移植至具備強力 GPU 或專屬 NPU 的計算平台（如 NVIDIA Jetson 系列）。透過平行化加速傅立葉轉換與矩陣運算，預期能將影像重建時間縮短至秒級，達成準即時的超解析度觀測，大幅降低研究人員的等待成本。

\subsection{多通道螢光同步成像}

本系統目前的濾光片架構偏向手動單色切換。未來規劃引入由步進馬達驅動的自動濾光片輪（Automatic Filter Wheel），並配合軟體腳本，實現 DAPI、GFP、RFP 等多種螢光染料的自動化、多通道同步掃描。結合大行程龍門載台，本機台將能完整支援 96 孔盤乃至 384 孔盤的高通量自動化藥物篩選實驗，充分發揮低成本高通量顯微鏡的核心價值。
